\chapter{2005年天津卷（理）}

\section{单选题}

\begin{problem}
设集合 \(A=\{x\mid |4x-1|\geqslant9,\ x\in\R\}\)，
\(B=\left\{x\mid \dfrac{x}{x+3}\geqslant0,\ x\in\R\right\}\)，则
\(A\cap B=\)（\(\quad\)）

\choices
    {\((-3,-2]\)}
    {\((-3,-2]\cup\left[0,\dfrac52\right]\)}
    {\((-\infty,-3]\cup\left[\dfrac52,+\infty\right)\)}
    {\((-\infty,-3)\cup\left[\dfrac52,+\infty\right)\)}
\end{problem}

\begin{answer}
D
\end{answer}

\begin{solution}
由 \(|4x-1|\geqslant9\)，得 \(4x-1\geqslant9\) 或 \(4x-1\leqslant-9\)，所以
\(A=(-\infty,-2]\cup[\frac52,+\infty)\). 对于分式不等式，定义域要求 \(x\ne-3\)；在
\((-\infty,-3)\)、\((-3,0)\) 和 \([0,+\infty)\) 上分别判断分子、分母的符号，得到
\(B=(-\infty,-3)\cup[0,+\infty)\). 因而
\(A\cap B=(-\infty,-3)\cup[\frac52,+\infty)\)，故选 D.
\end{solution}

\begin{problem}
若复数 \(\dfrac{a+3\i}{1+2\i}\)（\(a\in\R\)，\(\i\) 为虚数单位）是纯虚数，则实数 \(a\) 的值为（\(\quad\)）

\choices
    {\(-2\)}
    {\(4\)}
    {\(-6\)}
    {\(6\)}
\end{problem}

\begin{answer}
C
\end{answer}

\begin{solution}
分母有理化得
\[
\frac{a+3\i}{1+2\i}
=\frac{(a+3\i)(1-2\i)}5
=\frac{a+6}{5}+\frac{3-2a}{5}\i
\]
纯虚数的实部为零，所以 \(\dfrac{a+6}{5}=0\)，即 \(a=-6\)；此时虚部为
\(\dfrac{3-2a}{5}=3\ne0\)，确为纯虚数，故选 C.
\end{solution}

\begin{problem}
给出下列三个命题：

\begin{circlelist}
\item 若 \(a\geqslant b>-1\)，则 \(\dfrac{a}{1+a}\geqslant\dfrac{b}{1+b}\)；
\item 若正整数 \(m\) 和 \(n\) 满足 \(m\leqslant n\)，则 \(\sqrt{m(n-m)}\leqslant\dfrac n2\)；
\item 设 \(P(x_1,y_1)\) 为圆 \(O_1:x^2+y^2=9\) 上任一点，圆 \(O_2\) 以 \(Q(a,b)\) 为圆心且半径为 \(1\). 当 \((a-x_1)^2+(b-y_1)^2=1\) 时，圆 \(O_1\) 与圆 \(O_2\) 相切.
\end{circlelist}

其中假命题的个数为（\(\quad\)）

\choices
    {\(0\)}
    {\(1\)}
    {\(2\)}
    {\(3\)}
\end{problem}

\begin{answer}
B
\end{answer}

\begin{solution}
对命题（1），函数 \(g(x)=\dfrac{x}{1+x}\) 在 \(x>-1\) 上满足
\(g'(x)=\dfrac{1}{(1+x)^2}>0\)，所以 \(a\geqslant b\) 时
\(g(a)\geqslant g(b)\)，命题（1）正确.

对命题（2），两边非负，平方后只需证 \(4m(n-m)\leqslant n^2\). 而
\[
n^2-4m(n-m)=(n-2m)^2\geqslant0
\]
所以命题（2）正确. 命题（3）的条件只说明 \(P\) 同时在两圆上，不能推出两圆相切.
例如取 \(P=(3,0)\)，\(Q=(3,1)\)，则 \(P\) 在 \(O_1\) 上且
\(PQ=1\)，但两圆圆心距为 \(\sqrt{10}\)，满足
\(|3-1|<\sqrt{10}<3+1\)，两圆相交于两个点而不相切，所以命题（3）错误.
假命题有 \(1\) 个，故选 B.
\end{solution}

\begin{problem}
设 \(\alpha\)，\(\beta\)，\(\gamma\) 为平面，\(m\)，\(n\)，\(l\) 为直线，则
\(m\perp\beta\) 的一个充分条件是（\(\quad\)）

\choices
    {\(\alpha\perp\beta,\ \alpha\cap\beta=l,\ m\perp l\)}
    {\(\alpha\cap\gamma=m\)，\(\ \alpha\perp\gamma\)，\(\ \beta\perp\gamma\)}
    {\(\alpha\perp\gamma\)，\(\ \beta\perp\gamma\)，\(\ m\perp\alpha\)}
    {\(n\perp\alpha\)，\(\ n\perp\beta\)，\(\ m\perp\alpha\)}
\end{problem}

\begin{answer}
D
\end{answer}

\begin{solution}
选项 D 中，直线 \(n\) 同时垂直平面 \(\alpha\)、\(\beta\)，两平面的法向量都与
\(n\) 平行，因此 \(\alpha\parallel\beta\)（也允许两平面重合）. 又 \(m\perp\alpha\)，
所以 \(m\) 也垂直于与 \(\alpha\) 平行的平面 \(\beta\)，故 D 是充分条件. 选项 A
只给出 \(m\perp l\)，没有说明 \(m\) 在 \(\alpha\) 内. 例如取
\(\alpha:z=0\)、\(\beta:y=0\)、\(l\) 为 \(x\) 轴、\(m\) 为 \(z\) 轴，则 A 的条件成立，
但 \(m\subset\beta\)，所以 \(m\not\perp\beta\). 对于 B，取
\(\alpha=\beta:y=0\)、\(\gamma:z=0\)、\(m\) 为 \(x\) 轴，则 B 的条件成立，
但 \(m\subset\beta\). 对于 C，取 \(\alpha:y=0\)、\(\beta:x=0\)、\(\gamma:z=0\)、\(m\) 为 \(y\) 轴，
则 C 的条件成立，但 \(m\subset\beta\). 因此 A、B、C 均不能保证 \(m\perp\beta\).
\end{solution}

\begin{problem}
设双曲线以椭圆 \(\dfrac{x^2}{25}+\dfrac{y^2}{9}=1\) 长轴的两个端点为焦点，其准线过椭圆的焦点，则双曲线的渐近线的斜率为（\(\quad\)）

\choices
    {\(\pm2\)}
    {\(\pm\dfrac43\)}
    {\(\pm\dfrac12\)}
    {\(\pm\dfrac34\)}
\end{problem}

\begin{answer}
C
\end{answer}

\begin{solution}
椭圆长轴端点为 \((\pm5,0)\)，椭圆焦点为 \((\pm4,0)\). 设双曲线为
\(\dfrac{x^2}{p^2}-\dfrac{y^2}{q^2}=1\)，则焦半距 \(c=5\)，且准线
\(x=\pm\dfrac{p^2}{c}\) 经过 \((\pm4,0)\)，故 \(\dfrac{p^2}{5}=4\)，即
\(p^2=20\). 由 \(c^2=p^2+q^2\) 得 \(q^2=5\)，所以渐近线斜率为
\(\pm\dfrac qp=\pm\dfrac12\)，故选 C.
\end{solution}

\begin{problem}
从集合 \(\{1,2,3,\cdots,11\}\) 中任选两个元素作为椭圆方程
\(\dfrac{x^2}{m^2}+\dfrac{y^2}{n^2}=1\) 中的 \(m\) 和 \(n\)，则能组成落在矩形区域
\(B=\{(x,y)\mid |x|<11\text{ 且 }|y|<9\}\) 内的椭圆个数为（\(\quad\)）

\choices
    {\(43\)}
    {\(72\)}
    {\(86\)}
    {\(90\)}
\end{problem}

\begin{answer}
B
\end{answer}

\begin{solution}
椭圆上点的横坐标绝对值不超过 \(m\)，要使整个椭圆落在 \(|x|<11\) 内，必须且只需
\(m\leqslant10\)；同理，必须且只需 \(n\leqslant8\). 先把 \(m\) 和 \(n\) 看作有序取值，
共有 \(10\times8\) 种. 题目要求选取两个不同元素，因此要删去 \(m=n\) 且
\(m=n\in\{1,2,\ldots,8\}\) 的 \(8\) 种，故个数为
\(10\times8-8=72\)，选 B.
\end{solution}

\begin{problem}
某人射击一次击中的概率为 \(0.6\)，经过 \(3\) 次射击，此人至少有两次击中目标的概率为（\(\quad\)）

\choices
    {\(\dfrac{81}{125}\)}
    {\(\dfrac{54}{125}\)}
    {\(\dfrac{36}{125}\)}
    {\(\dfrac{27}{125}\)}
\end{problem}

\begin{answer}
A
\end{answer}

\begin{solution}
按独立重复射击计算，设 \(p=0.6=\dfrac35\)，则至少击中两次包括恰好两次和三次. 因此所求概率为
\[
\mathrm C_3^2p^2(1-p)+p^3
=3\left(\frac35\right)^2\frac25+\left(\frac35\right)^3
=\frac{54}{125}+\frac{27}{125}
=\frac{81}{125}
\]
故选 A.
\end{solution}

\begin{problem}
要得到函数 \(y=\sqrt2\cos x\) 的图象，只需将函数
\(y=\sqrt2\sin\left(2x+\dfrac\pi4\right)\) 的图象上所有的点的（\(\quad\)）

\choices
    {横坐标缩短到原来的 \(\dfrac12\) 倍（纵坐标不变），再向左平行移动 \(\dfrac\pi8\) 个单位长度}
    {横坐标缩短到原来的 \(\dfrac12\) 倍（纵坐标不变），再向右平行移动 \(\dfrac\pi4\) 个单位长度}
    {横坐标伸长到原来的 \(2\) 倍（纵坐标不变），再向左平行移动 \(\dfrac\pi4\) 个单位长度}
    {横坐标伸长到原来的 \(2\) 倍（纵坐标不变），再向右平行移动 \(\dfrac\pi8\) 个单位长度}
\end{problem}

\begin{answer}
C
\end{answer}

\begin{solution}
横坐标伸长为原来的 \(2\) 倍后，函数变为
\(y=\sqrt2\sin\left(x+\dfrac\pi4\right)\). 再向左平移 \(\dfrac\pi4\) 个单位，
函数变为 \(y=\sqrt2\sin\left(x+\dfrac\pi2\right)=\sqrt2\cos x\)，故选 C.
\end{solution}

\begin{problem}
设 \(f^{-1}(x)\) 是函数 \(f(x)=\dfrac12(a^x-a^{-x})\)（\(a>1\)）的反函数，
则使 \(f^{-1}(x)>1\) 成立的 \(x\) 的取值范围为（\(\quad\)）

\choices
    {\(\left(\dfrac{a^2-1}{2a},+\infty\right)\)}
    {\(\left(-\infty,\dfrac{a^2-1}{2a}\right)\)}
    {\(\left(\dfrac{a^2-1}{2a},a\right)\)}
    {\([a,+\infty)\)}
\end{problem}

\begin{answer}
A
\end{answer}

\begin{solution}
因为 \(a>1\)，有 \(f'(x)=\dfrac{\ln a}{2}(a^x+a^{-x})>0\)，
所以 \(f\) 严格递增，且其值域为 \(\R\). 因而
\(f^{-1}(x)>1\) 等价于 \(x>f(1)\)，而
\[
f(1)=\frac12(a-a^{-1})=\frac{a^2-1}{2a}
\]
所求范围为 \(\left(\dfrac{a^2-1}{2a},+\infty\right)\)，故选 A.
\end{solution}

\begin{problem}
若函数 \(f(x)=\log_a(x^3-ax)\)（\(a>0\)，\(a\ne1\)）在区间
\(\left(-\dfrac12,0\right)\) 内单调递增，则 \(a\) 的取值范围是（\(\quad\)）

\choices
    {\(\left[\dfrac14,1\right)\)}
    {\(\left[\dfrac34,1\right)\)}
    {\(\left(\dfrac94,+\infty\right)\)}
    {\(\left(1,\dfrac94\right)\)}
\end{problem}

\begin{answer}
B
\end{answer}

\begin{solution}
在 \(x\in\left(-\dfrac12,0\right)\) 时 \(x<0\)，所以真数
\(x^3-ax=x(x^2-a)>0\) 等价于 \(a>x^2\). 要使真数在整个开区间内有定义，
必要充分条件为 \(a\geqslant\dfrac14\).

令 \(g(x)=x^3-ax\)，则 \(g'(x)=3x^2-a\). 若 \(0<a<1\)，对数函数递减，
要使 \(\log_a g(x)\) 递增，需在整个区间内有 \(g'(x)<0\). 当
\(a\geqslant\dfrac34\) 时，由 \(x^2<\dfrac14\) 得
\(g'(x)<\dfrac34-a\leqslant0\)，且不等号实际为严格不等号；当
\(\dfrac14\leqslant a<\dfrac34\) 时，可取充分接近 \(-\dfrac12\) 的 \(x\)，使
\(3x^2-a>0\)，此时复合函数的导数为负，不能在整个区间内递增. 若 \(a>1\)，
则需 \(g'(x)>0\)，但 \(g'(x)<\dfrac34-a<0\)，不可能. 结合 \(a\ne1\)，得
\(a\in\left[\dfrac34,1\right)\)，故选 B.
\end{solution}

\section{填空题}

\begin{problem}
设 \(n\in\mathbb N^*\)，则
\(\mathrm C_n^1+\mathrm C_n^2 6+\mathrm C_n^3 6^2+\cdots+\mathrm C_n^n6^{n-1}=\)\fillinblank{}.
\end{problem}

\begin{answer}
\(\dfrac{7^n-1}{6}\)
\end{answer}

\begin{solution}
由二项式定理，
\[
7^n=(1+6)^n=1+\mathrm C_n^1 6+\mathrm C_n^2 6^2+\cdots+\mathrm C_n^n6^n
\]
两边减去 \(1\) 后除以 \(6\)，得到原式为 \(\dfrac{7^n-1}{6}\).
\end{solution}

\begin{problem}
如图，\(PA\perp\) 平面 \(ABC\)，\(\angle ACB=90^\circ\) 且
\(PA=AC=BC=a\)，则异面直线 \(PB\) 与 \(AC\) 所成角的正切值等于\fillinblank{}.

\bitmapfigure[max width=0.28\linewidth]{img/2005/tianjin_science/q12_fig1.png}
\end{problem}

\begin{answer}
\(\sqrt2\)
\end{answer}

\begin{solution}
建立直角坐标系，取 \(A=(0,0,0)\)，\(C=(a,0,0)\)，\(B=(a,a,0)\)，\(P=(0,0,a)\).
则 \(\overrightarrow{AC}=(a,0,0)\)，\(\overrightarrow{PB}=(a,a,-a)\). 设两直线所成的锐角为
\(\theta\)，则
\[
\cos\theta
=\frac{|\overrightarrow{AC}\cdot\overrightarrow{PB}|}{|AC|\,|PB|}
=\frac{a^2}{a\cdot\sqrt3a}
=\frac1{\sqrt3}
\]
因而 \(\tan\theta=\sqrt{\dfrac1{\cos^2\theta}-1}=\sqrt2\).
\end{solution}

\begin{problem}
在数列 \(\{a_n\}\) 中，\(a_1=1\)，\(a_2=2\)，且
\(a_{n+2}-a_n=1+(-1)^n\)（\(n\in\mathbb N^*\)）. 则
\(S_{100}=\)\fillinblank{}.
\end{problem}

\begin{answer}
\(2600\)
\end{answer}

\begin{solution}
当 \(n\) 为奇数时，\(a_{n+2}=a_n\)，所以所有奇数项均为 \(1\). 当 \(n=2k\) 为偶数时，
\(a_{2k+2}=a_{2k}+2\)，结合 \(a_2=2\) 得 \(a_{2k}=2k\). 因而
\[
S_{100}
=\sum_{k=1}^{50}a_{2k-1}+\sum_{k=1}^{50}a_{2k}
=50+2\sum_{k=1}^{50}k
=50+2550
=2600
\]
\end{solution}

\begin{problem}
在直角坐标系 \(xOy\) 中，已知点 \(A(0,1)\) 和点 \(B(-3,4)\)，若点 \(C\) 在
\(\angle AOB\) 的平分线上且 \(\left|\overrightarrow{OC}\right|=2\)，则
\(\overrightarrow{OC}=\)\fillinblank{}.
\end{problem}

\begin{answer}
\(\left(-\dfrac{\sqrt{10}}5,\dfrac{3\sqrt{10}}5\right)\)
\end{answer}

\begin{solution}
\(\overrightarrow{OA}=(0,1)\)，\(\overrightarrow{OB}=(-3,4)\)，且 \(|OA|=1\)，\(|OB|=5\).
内角平分线方向可取单位向量之和
\[
\frac{\overrightarrow{OA}}{|OA|}+\frac{\overrightarrow{OB}}{|OB|}
=\left(-\frac35,\frac95\right)
\]
其方向为 \((-1,3)\). 该方向的单位向量为
\(\dfrac1{\sqrt{10}}(-1,3)\)，再乘以 \(|OC|=2\)，得
\[
\overrightarrow{OC}
=\left(-\frac2{\sqrt{10}},\frac6{\sqrt{10}}\right)
=\left(-\frac{\sqrt{10}}5,\frac{3\sqrt{10}}5\right)
\]
\end{solution}

\begin{problem}
某公司有 \(5\text{ 万元}\) 资金用于投资开发项目，如果成功，一年后可获利 \(12\%\)，
一旦失败，一年后将丧失全部资金的 \(50\%\). 下表是过去 \(200\) 例类似项目开发的实施结果：

\begin{center}
\begin{tabular}{|c|c|}
\hline
投资成功 & 投资失败 \\
\hline
\(192\) 次 & \(8\) 次 \\
\hline
\end{tabular}
\end{center}

则该公司一年后估计可获收益的期望是\fillinblank{}\(\text{元}\).
\end{problem}

\begin{answer}
\(4760\)
\end{answer}

\begin{solution}
由历史数据估计成功、失败的概率分别为
\(\dfrac{192}{200}=0.96\) 和 \(\dfrac8{200}=0.04\). 成功时收益为
\(50000\times12\%=6000\) 元，失败时收益为
\(-50000\times50\%=-25000\) 元. 所以收益的期望为
\[
0.96\times6000+0.04\times(-25000)=5760-1000=4760\text{ 元}
\]
\end{solution}

\begin{problem}
设 \(f(x)\) 是定义在 \(\R\) 上的奇函数，且 \(y=f(x)\) 的图象关于直线
\(x=\dfrac12\) 对称，则
\(f(1)+f(2)+f(3)+f(4)+f(5)=\)\fillinblank{}.
\end{problem}

\begin{answer}
\(0\)
\end{answer}

\begin{solution}
关于直线 \(x=\dfrac12\) 对称意味着 \(f(x)=f(1-x)\). 取 \(x=1\)，得
\(f(1)=f(0)=0\). 再由
\[
f(n)=f(1-n)=-f(n-1)
\]
并结合 \(f(1)=0\) 归纳得 \(f(2)=f(3)=f(4)=f(5)=0\)，故所求和为 \(0\).
\end{solution}

\section{解答题}

\begin{problem}
在 \(\triangle ABC\) 中，\(\angle A\)，\(\angle B\)，\(\angle C\) 所对的边长分别为
\(a\)，\(b\)，\(c\). 设 \(a\)，\(b\)，\(c\) 满足条件
\(b^2+c^2-bc=a^2\) 和 \(\dfrac cb=\dfrac12+\sqrt3\)，求
\(\angle A\) 和 \(\tan B\) 的值.
\end{problem}

\begin{solution}
由余弦定理，\(a^2=b^2+c^2-2bc\cos A\). 与已知等式比较得
\(2bc\cos A=bc\)，所以 \(\cos A=\dfrac12\)，从而 \(A=60^\circ\).
又 \(C=120^\circ-B\)，由正弦定理
\[
\frac cb
=\frac{\sin C}{\sin B}
=\frac{\sin(120^\circ-B)}{\sin B}
=\frac{\sqrt3}{2}\cot B+\frac12
\]
代入 \(\dfrac cb=\dfrac12+\sqrt3\)，得 \(\dfrac{\sqrt3}{2}\cot B=\sqrt3\)，即
\(\cot B=2\)，所以 \(\tan B=\dfrac12\).
\end{solution}

\begin{problem}
已知 \(u_n=a^n+a^{n-1}b+a^{n-2}b^2+\cdots+ab^{n-1}+b^n\)
\((n\in\mathbb N^*,\ a>0,\ b>0)\).

\begin{enumerate}
\item 当 \(a=b\) 时，求数列 \(\{u_n\}\) 的前 \(n\) 项和 \(S_n\)；
\item 求 \(\displaystyle\lim_{n\to\infty}\frac{u_n}{u_{n-1}}\).
\end{enumerate}
\end{problem}

\begin{solution}
\begin{enumerate}
\item 当 \(a=b\) 时，原式的 \(k+1\) 项均为 \(a^k\)，所以 \(u_k=(k+1)a^k\). 当
\(a\ne1\) 时，分别使用加权等比数列求和公式和等比数列求和公式，得
\[
\begin{gathered}
S_n=\sum_{k=1}^n(k+1)a^k=\sum_{k=1}^nka^k+\sum_{k=1}^na^k,\\
S_n=\frac{a-(n+1)a^{n+1}+na^{n+2}}{(1-a)^2}
  +\frac{a(1-a^n)}{1-a},\qquad a\ne1.
\end{gathered}
\]
当 \(a=1\) 时，\(u_k=k+1\)，所以
\[
S_n=\sum_{k=1}^n(k+1)=\frac{n(n+3)}2
\]
\item 当 \(a\ne b\) 时，等比数列求和得
\[
u_n=\frac{a^{n+1}-b^{n+1}}{a-b}
\]
若 \(a>b\)，则
\[
\frac{u_n}{u_{n-1}}
=a\frac{1-(b/a)^{n+1}}{1-(b/a)^n}\longrightarrow a
\]
所以此时极限为 \(a\).
若 \(b>a\)，令 \(s=\dfrac ab\in(0,1)\)，则
\[
\frac{u_n}{u_{n-1}}
=b\frac{1-s^{n+1}}{1-s^n}\longrightarrow b
\]
当 \(a=b\) 时，
\[
\frac{u_n}{u_{n-1}}=a\frac{n+1}{n}\longrightarrow a
\]
综合三种情形，极限为 \(\max\{a,b\}\).
\end{enumerate}
\end{solution}

\begin{problem}
如图，在斜三棱柱 \(ABC-A_1B_1C_1\) 中，
\(\angle A_1AB=\angle A_1AC\)，\(AB=AC\)，\(A_1A=A_1B=a\)，侧面
\(B_1BCC_1\) 与底面 \(ABC\) 所成的二面角为 \(120^\circ\)，\(E\)，\(F\) 分别是棱
\(B_1C_1\)，\(A_1A\) 的中点.

\begin{enumerate}
\item 求 \(A_1A\) 与底面 \(ABC\) 所成的角；
\item 证明 \(A_1E\parallel\text{平面 }B_1FC\)；
\item 求经过 \(A_1\)，\(A\)，\(B\)，\(C\) 四点的球的体积.
\end{enumerate}

\bitmapfigure[max width=0.42\linewidth]{img/2005/tianjin_science/q19_fig1.png}
\end{problem}

\begin{solution}
\begin{enumerate}
\item 设 \(D\) 为棱 \(BC\) 的中点，\(H\) 为点 \(A_1\) 在底面 \(ABC\) 内的射影. 因
\(A_1H\perp\) 底面，且 \(AB=AC\)，由
\(\angle A_1AB=\angle A_1AC\) 得
\(\overrightarrow{AH}\cdot\overrightarrow{AB}=
\overrightarrow{AH}\cdot\overrightarrow{AC}\)，所以
\(AH\perp BC\). 又 \(AD\perp BC\)，故 \(A\)，\(H\)，\(D\) 共线. 取底面为
\(z=0\)，以 \(AD\) 为 \(x\) 轴并令 \(D\) 的横坐标为正，可取
\(A=(0,0,0)\)，\(A_1=(p,0,h)\)，\(B=(q,t,0)\)，\(C=(q,-t,0)\)
其中 \(q>0\)，\(t>0\)，且 \(h>0\). 过 \(D\) 作垂直于 \(BC\) 的截面，底面内的
截线方向为 \((-1,0,0)\)，侧面内的截线方向为 \((p,0,h)\). 因二面角为
\(120^\circ\)，有
\[
\cos120^\circ=\frac{(-1,0,0)\cdot(p,0,h)}{\sqrt{p^2+h^2}}=-\frac p{AA_1}
\]
又 \(AA_1=a\)，所以 \(p=\dfrac a2\)，再由
\[
h=\sqrt{a^2-p^2}=\frac{\sqrt3}{2}a
\]
因而 \(A_1A\) 与底面所成的角 \(\theta\) 满足
\(\sin\theta=\dfrac ha=\dfrac{\sqrt3}{2}\)，故 \(\theta=60^\circ\).
\item 由平移关系，
\(B_1=(q+p,t,h)\)，\(C_1=(q+p,-t,h)\)
所以
\(E=(q+p,0,h)\)，\(F=\left(\frac p2,0,\frac h2\right)\).
因此 \(\overrightarrow{A_1E}=(q,0,0)\). 另一方面，
\(\overrightarrow{FB_1}=\left(q+\frac p2,t,\frac h2\right)\)，\(\overrightarrow{FC}=\left(q-\frac p2,-t,-\frac h2\right)\)
从而 \(\overrightarrow{FB_1}+\overrightarrow{FC}=(2q,0,0)\). 这两个向量都在
平面 \(B_1FC\) 内，且 \(q>0\)，故该平面内存在一条与 \(x\) 轴平行的直线.
又由 \(A_1B=a\) 及 \(p^2+h^2=a^2\) 得
\(q^2+t^2=2pq=aq\). 而 \(\overrightarrow{A_1E}=(q,0,0)\) 也与 \(x\) 轴平行.
\(\overrightarrow{FB_1}-\overrightarrow{FC}=(p,2t,h)\)，故平面 \(B_1FC\) 内过 \(F\) 的任一向量可写成
\(r(1,0,0)+s(p,2t,h)\)；若它等于
\(\overrightarrow{FA_1}=(p/2,0,h/2)\)，由第二坐标得 \(s=0\)，但第三坐标随即不可能相等.
所以 \(A_1\notin\text{平面 }B_1FC\)，故 \(A_1E\parallel\text{平面 }B_1FC\).
\item 由 \(A_1B=a\) 以及 \(p=\dfrac a2\)，\(h=\dfrac{\sqrt3}2a\)，得
\((q-p)^2+t^2+h^2=a^2\)，\(\Longrightarrow\)，\(q^2+t^2=aq\).
所以 \(\triangle ABC\) 的外接圆圆心为 \(O=(\frac a2,0,0)\). 设球心为
\(S=(\frac a2,0,z)\)，由 \(SA=SA_1\) 得
\[
\frac{a^2}{4}+z^2=(z-h)^2
\]
故 \(z=\dfrac a{2\sqrt3}\). 于是球半径满足
\[
R^2=SA^2=\frac{a^2}{4}+\frac{a^2}{12}=\frac{a^2}{3}
\]
所以
\[
V=\frac43\pi R^3=\frac{4\sqrt3}{27}\pi a^3
\]
\end{enumerate}
\end{solution}

\begin{problem}
某人在一山坡 \(P\) 处观看对面山顶上的一座铁塔，如图所示，塔高
\(BC=80\text{ 米}\)，塔所在的山高 \(OB=220\text{ 米}\)，\(OA=200\text{ 米}\).
图中所示的山坡可视为直线 \(l\) 且点 \(P\) 在直线 \(l\) 上，\(l\) 与水平地面的夹角为
\(\alpha\)，\(\tan\alpha=\dfrac12\). 试问此人距水平地面多高时，观看塔的视角
\(\angle BPC\) 最大？（不计此人的身高）

\bitmapfigure[max width=0.56\linewidth]{img/2005/tianjin_science/q20_fig1.png}
\end{problem}

\begin{solution}
设 \(P\) 距水平地面的高度为 \(h\geqslant0\). 取 \(O\) 为原点、水平向右为 \(x\) 轴，
则 \(A=(200,0)\)，\(B=(0,220)\)，\(C=(0,300)\). 由于山坡斜率为
\(\dfrac12\)，有 \(P=(200+2h,h)\). 记 \(x=200+2h\)，则
\(\overrightarrow{PB}=(-x,220-h)\)，\(\overrightarrow{PC}=(-x,300-h)\).
视角为锐角，且
\[
\tan\angle BPC
=\frac{|\overrightarrow{PB}\times\overrightarrow{PC}|}
{\overrightarrow{PB}\cdot\overrightarrow{PC}}
=\frac{80x}{x^2+(220-h)(300-h)}
=\frac{160(h+100)}{5h^2+280h+106000}
\]
令右端为 \(T(h)\)，则 \(T'(h)\) 的符号由
\[
5h^2+280h+106000-(h+100)(10h+280)
=-5(h-60)(h+260)
\]
决定. 所以 \(T(h)\) 在 \(0\leqslant h<60\) 上递增，在 \(h>60\) 上递减，
最大值在 \(h=60\) 时取得. 因而此人距水平地面 \(60\) 米高时视角最大.
\end{solution}

\begin{problem}
抛物线 \(C\) 的方程为 \(y=ax^2\)（\(a<0\)），过抛物线 \(C\) 上一点
\(P(x_0,y_0)\)（\(x_0\ne0\)）作斜率为 \(k_1\)，\(k_2\) 的两条直线分别交抛物线
\(C\) 于 \(A(x_1,y_1)\)，\(B(x_2,y_2)\) 两点（\(P\)，\(A\)，\(B\) 三点互不相同），
且满足 \(k_2+\lambda k_1=0\)（\(\lambda\ne0\)，\(\lambda\ne-1\)）.

\begin{enumerate}
\item 求抛物线 \(C\) 的焦点坐标和准线方程；
\item 设直线 \(AB\) 上一点 \(M\)，满足 \(\overrightarrow{BM}=\lambda\overrightarrow{MA}\)，
证明线段 \(PM\) 的中点在 \(y\) 轴上；
\item 当 \(\lambda=1\) 时，若点 \(P\) 的坐标为 \((1,-1)\)，求
\(\angle PAB\) 为钝角时点 \(A\) 的纵坐标 \(y_1\) 的取值范围.
\end{enumerate}
\end{problem}

\begin{solution}
\begin{enumerate}
\item 将 \(y=ax^2\) 写成 \(x^2=\dfrac1a y=4py\)，得 \(p=\dfrac1{4a}\).
所以焦点为 \(\left(0,\dfrac1{4a}\right)\)，准线为 \(y=-\dfrac1{4a}\).
\item 抛物线上两点的弦斜率满足
\[
\frac{y_1-y_0}{x_1-x_0}=a(x_1+x_0)
\]
故 \(k_1=a(x_0+x_1)\)，\(k_2=a(x_0+x_2)\). 由条件得
\[
x_2+\lambda x_1+(1+\lambda)x_0=0
\]
又 \(\overrightarrow{BM}=\lambda\overrightarrow{MA}\) 给出
\((1+\lambda)M=B+\lambda A\)，所以
\[
x_M=\frac{x_2+\lambda x_1}{1+\lambda}=-x_0
\]
因而 \(PM\) 的中点横坐标为 \(\dfrac{x_0-x_0}{2}=0\)，即线段 \(PM\) 的中点在
\(y\) 轴上.
\item \(P=(1,-1)\) 在抛物线上，故 \(a=-1\)，抛物线为 \(y=-x^2\). 令 \(x_1=t\)，
由 \(k_1+k_2=0\) 得 \(x_2=-2-t\). 于是
\(A=(t,-t^2)\)，\(B=(-2-t,-(t+2)^2)\).
并且
\(\overrightarrow{AP}=(1-t)(1,-t-1)\)，\(\overrightarrow{AB}=-2(t+1)(1,2)\).
因而
\[
\overrightarrow{AP}\cdot\overrightarrow{AB}
=2(1-t)(t+1)(2t+1)
\]
\(\angle PAB\) 为钝角等价于该内积小于零，即
\[
t\in\left(-1,-\frac12\right)\cup(1,+\infty)
\]
题设还排除 \(t=-1\) 时 \(A=B\)、\(t=1\) 时 \(A=P\) 以及 \(t=-3\) 时
\(B=P\) 的情形；其中 \(t=-3\) 不在上面的钝角范围内. 因为 \(y_1=-t^2\)，故
\[
y_1\in(-\infty,-1)\cup\left(-1,-\frac14\right)
\]
\end{enumerate}
\end{solution}

\begin{problem}
设函数 \(f(x)=x\sin x\)（\(x\in\R\)）.

\begin{enumerate}
\item 证明 \(f(x+2k\pi)-f(x)=2k\pi\sin x\)，其中 \(k\) 为整数；
\item 设 \(x_0\) 为 \(f(x)\) 的一个极值点，证明
\(\left[f(x_0)\right]^2=\dfrac{x_0^4}{1+x_0^2}\)；
\item 设 \(f(x)\) 在 \((0,+\infty)\) 内的全部极值点按从小到大的顺序排列
\(a_1\)，\(a_2\)，\(\ldots\)，\(a_n\)，\(\ldots\)，证明
\(\dfrac\pi2<a_{n+1}-a_n<\pi\)（\(n=1\)，\(2\)，\(\ldots\)）.
\end{enumerate}
\end{problem}

\begin{solution}
\begin{enumerate}
\item 因为 \(\sin(x+2k\pi)=\sin x\)，所以
\[
f(x+2k\pi)-f(x)
=(x+2k\pi)\sin x-x\sin x
=2k\pi\sin x
\]
\item \(f'(x)=\sin x+x\cos x\). 在极值点 \(x_0\) 有
\(\sin x_0+x_0\cos x_0=0\)，即 \(\sin x_0=-x_0\cos x_0\). 与
\(\sin^2x_0+\cos^2x_0=1\) 联立，得
\[
(1+x_0^2)\cos^2x_0=1
\]
所以 \(\sin^2x_0=\dfrac{x_0^2}{1+x_0^2}\).
因此
\[
\left[f(x_0)\right]^2=x_0^2\sin^2x_0
=\frac{x_0^4}{1+x_0^2}
\]
\item 极值点不能满足 \(\cos x=0\)，因为此时
\(f'(x)=\sin x\ne0\). 因而正极值点必须满足 \(\tan x=-x\). 对 \(x>0\)，
右端 \(-x<0\)，所以只需考察正切函数为负的区间
\(I_n=\left(\left(n-\frac12\right)\pi,n\pi\right)\)，\(n=1\)，\(2\)，\(\ldots\)
令 \(g(x)=\tan x+x\). 在每个 \(I_n\) 上，
\(g'(x)=\sec^2x+1>0\)，且当 \(x\) 从区间左端趋近时 \(g(x)\to-\infty\)，而
\(g(n\pi)=n\pi>0\)，所以每个 \(I_n\) 内恰有一个根，记为 \(a_n\). 根的两侧
\(g\) 符号相反，而 \(\cos x\) 在 \(I_n\) 内不变号，因此
\(f'(x)=\cos x\,g(x)\) 在 \(a_n\) 两侧变号；这些根恰为全部正极值点. 由
\(a_n<n\pi\) 及 \(a_{n+1}>\left(n+\frac12\right)\pi\)，得
\[
a_n+\frac\pi2<\left(n+\frac12\right)\pi<a_{n+1}
\]
所以 \(a_{n+1}-a_n>\dfrac\pi2\). 又因为
\(a_n>\left(n-\frac12\right)\pi\) 且 \(a_n<n\pi\)，有 \(a_n+\pi\in I_{n+1}\)，且
\[
g(a_n+\pi)=\tan a_n+a_n+\pi=\pi>0
\]
由于 \(g\) 在 \(I_{n+1}\) 上递增且在根处为零，根 \(a_{n+1}\) 位于 \(a_n+\pi\) 左侧，
故 \(a_{n+1}-a_n<\pi\). 结论成立.
\end{enumerate}
\end{solution}
